% ============================================================
% SEÇÃO 9 – CASOS DE TESTE
% ============================================================

\section{CASOS DE TESTE}

\subsection*{Caso de teste \#1 – Cadastro do usuário}

\textbf{Propósito:} Verificar se o usuário consegue criar uma conta no LifeCity
informando dados válidos e se o sistema apresenta mensagens adequadas em
casos de erro.

\textbf{Pré-condições:}
\begin{itemize}
  \item Aplicativo instalado e com acesso à internet.
  \item Usuário ainda não autenticado.
\end{itemize}

\textbf{Saída aguardada:}
\begin{itemize}
  \item Cadastro realizado com sucesso e redirecionamento para tela de login;
  \item Em caso de erro, exibir mensagem explicando o problema e não concluir o cadastro.
\end{itemize}

\begin{longtable}{|c|p{1.5cm}|p{2cm}|p{2.5cm}|p{2cm}|p{2cm}|p{2.5cm}|}
  \hline
  \textbf{\#} & \textbf{Nome} & \textbf{CPF} & \textbf{E-mail} & \textbf{Telefone} & \textbf{Senha} & \textbf{Saída esperada} \\
  \hline
  \endfirsthead
  \hline
  \textbf{\#} & \textbf{Nome} & \textbf{CPF} & \textbf{E-mail} & \textbf{Telefone} & \textbf{Senha} & \textbf{Saída esperada} \\
  \hline
  \endhead
  1 & Ana Silva & 123.456.789-09 & ana@email.com & (19) 99999-0000 & Senha@123 & Cadastro realizado com sucesso. \\
  \hline
  2 & (vazio) & 123.456.789-09 & ana@email.com & (19) 99999-0000 & Senha@123 & ``Preencha o campo Nome.'' \\
  \hline
  3 & João & 123.456.789-0 & joao@email.com & (19) 98888-0000 & Senha@123 & ``CPF inválido.'' \\
  \hline
  4 & João & 123.456.789-09 & joao@ & (19) 98888-0000 & Senha@123 & ``E-mail inválido.'' \\
  \hline
  5 & João & 123.456.789-09 & joao@email.com & 9888-000 & Senha@123 & ``Telefone inválido.'' \\
  \hline
  6 & João & 123.456.789-09 & joao@email.com & (19) 98888-0000 & 123 & ``Senha deve atender aos requisitos mínimos.'' \\
  \hline
  7 & João & 123.456.789-09 & joao@email.com & (19) 98888-0000 & Senha@123 / SenhaErrada & ``As senhas informadas não conferem.'' \\
  \hline
  8 & João & (vazio) & (vazio) & (vazio) & (vazio) & Mensagens de campos obrigatórios para CPF, E-mail, Telefone e Senha. \\
  \hline
\end{longtable}

\subsection*{Caso de teste \#2 – Login e visualização do mapa}

\textbf{Propósito:} Validar se o usuário consegue autenticar-se no sistema e
visualizar o mapa com eventos/ocorrências da cidade.

\textbf{Pré-condições:}
\begin{itemize}
  \item Usuário já cadastrado no sistema.
  \item Aplicativo com acesso à internet.
\end{itemize}

\begin{longtable}{|c|p{4cm}|p{2.5cm}|p{5.5cm}|}
  \hline
  \textbf{\#} & \textbf{E-mail} & \textbf{Senha} & \textbf{Saída esperada} \\
  \hline
  \endfirsthead
  \hline
  \textbf{\#} & \textbf{E-mail} & \textbf{Senha} & \textbf{Saída esperada} \\
  \hline
  \endhead
  1 & ana@email.com & Senha@123 & Login realizado; usuário redirecionado para tela principal com mapa carregado. \\
  \hline
  2 & naoexiste@email.com & Senha@123 & ``Usuário não encontrado ou credenciais inválidas.'' \\
  \hline
  3 & (vazio) & Senha@123 & ``Campo de e-mail obrigatório.'' \\
  \hline
  4 & ana@email.com & (vazio) & ``Campo de senha obrigatório.'' \\
  \hline
  5 & ana@email.com & senhaErrada & ``Usuário não encontrado ou credenciais inválidas''; não realiza login. \\
  \hline
\end{longtable}

\subsection*{Caso de teste \#3 – Registro de reclamação/evento}

\textbf{Propósito:} Verificar se o usuário consegue publicar reclamações e
eventos e se o sistema exibe mensagens adequadas para campos inválidos.

\begin{longtable}{|c|p{2cm}|p{3cm}|p{2cm}|p{2cm}|p{3.5cm}|}
  \hline
  \textbf{\#} & \textbf{Tipo} & \textbf{Descrição} & \textbf{Categ.} & \textbf{Local.} & \textbf{Saída esperada} \\
  \hline
  \endfirsthead
  \hline
  \textbf{\#} & \textbf{Tipo} & \textbf{Descrição} & \textbf{Categ.} & \textbf{Local.} & \textbf{Saída esperada} \\
  \hline
  \endhead
  1 & Reclamação & Buraco na Rua X & Infraestrutura & GPS ok & Reclamação registrada com sucesso e exibida no mapa. \\
  \hline
  2 & Evento & Feira comunitária & Evento & Praça Y & Evento registrado com sucesso e exibido no mapa. \\
  \hline
  3 & Reclamação & (vazio) & Limpeza & GPS ok & ``Preencha a descrição da publicação.''; não registra. \\
  \hline
  4 & Evento & Aulão de esportes & Esportes & GPS ok & ``Informe data e horário do evento.''; não registra. \\
  \hline
  5 & Reclamação & Luz queimada & Infraestrutura & GPS negado & Sistema oferece campo de endereço manual. \\
  \hline
\end{longtable}

\subsection*{Caso de teste \#4 – Filtrar publicações por categoria}

\textbf{Propósito:} Verificar se o usuário consegue aplicar filtros de categoria
no mapa e se o sistema exibe apenas as publicações compatíveis com a
categoria selecionada.

\begin{longtable}{|c|p{3cm}|p{4cm}|p{4.5cm}|}
  \hline
  \textbf{\#} & \textbf{Categoria selecionada} & \textbf{Situação na base} & \textbf{Saída esperada} \\
  \hline
  \endfirsthead
  \hline
  \textbf{\#} & \textbf{Categoria selecionada} & \textbf{Situação na base} & \textbf{Saída esperada} \\
  \hline
  \endhead
  1 & Infraestrutura & Publicações de Infraestrutura e Limpeza & Exibir apenas publicações de Infraestrutura no mapa. \\
  \hline
  2 & Segurança & Publicações de Segurança e Esportes & Exibir apenas publicações de Segurança no mapa. \\
  \hline
  3 & Limpeza & Nenhuma publicação de Limpeza & Mapa sem marcadores. \\
  \hline
  4 & Esportes & Publicações de todas as categorias & Exibir apenas publicações de Esportes. \\
  \hline
  5 & (Todas / sem filtro) & Publicações de diferentes categorias & Exibir todas as publicações no mapa sem filtragem. \\
  \hline
\end{longtable}

\subsection*{Caso de teste \#5 – Desbloqueio de conquista}

\textbf{Propósito:} Verificar se o sistema desbloqueia conquistas automaticamente
ao atingir marcos predefinidos e emite a notificação correspondente.

\textbf{Pré-condições:}
\begin{itemize}
  \item Usuário logado, sem nenhuma conquista desbloqueada.
\end{itemize}

\begin{longtable}{|c|p{4cm}|p{3.5cm}|p{4.5cm}|}
  \hline
  \textbf{\#} & \textbf{Ação} & \textbf{Estado esperado} & \textbf{Saída esperada} \\
  \hline
  \endfirsthead
  \hline
  \textbf{\#} & \textbf{Ação} & \textbf{Estado esperado} & \textbf{Saída esperada} \\
  \hline
  \endhead
  1 & Usuário cria sua primeira reclamação & Conquista ``Primeira Voz'' ainda não desbloqueada & Conquista desbloqueada; notificação \texttt{achievement\_unlocked} emitida; 25 XP adicionados. \\
  \hline
  2 & Usuário cria a primeira reclamação uma segunda vez (mesmo usuário) & Conquista ``Primeira Voz'' já desbloqueada & Conquista não duplicada; sem nova notificação. \\
  \hline
  3 & Usuário recebe 10 curtidas em reclamações & Conquista ``Popular'' ainda não desbloqueada & Conquista desbloqueada; notificação emitida; 50 XP adicionados. \\
  \hline
\end{longtable}

\subsection*{Caso de teste \#6 – Criar equipe e convidar membro}

\textbf{Propósito:} Verificar o fluxo completo de criação de equipe e gerenciamento de membros.

\textbf{Pré-condições:}
\begin{itemize}
  \item Dois usuários cadastrados (criador e convidado).
  \item Usuário criador logado.
\end{itemize}

\begin{longtable}{|c|p{4cm}|p{2.5cm}|p{5cm}|}
  \hline
  \textbf{\#} & \textbf{Ação} & \textbf{Entrada} & \textbf{Saída esperada} \\
  \hline
  \endfirsthead
  \hline
  \textbf{\#} & \textbf{Ação} & \textbf{Entrada} & \textbf{Saída esperada} \\
  \hline
  \endhead
  1 & Criar equipe com nome válido & ``Guardiões do Bairro'' & Equipe criada; criador redirecionado para tela de detalhes. \\
  \hline
  2 & Criar equipe com nome vazio & (vazio) & ``Informe um nome para a equipe.''; não cria. \\
  \hline
  3 & Convidar usuário amigo para a equipe & ID do convidado & Convite enviado; convidado recebe notificação \texttt{team\_invite}. \\
  \hline
  4 & Convidado aceita o convite & — & Status do membro muda para \texttt{active}; equipe com 2 membros ativos. \\
  \hline
  5 & Convidado recusa o convite & — & Status do membro muda para \texttt{rejected}; equipe permanece com 1 membro ativo. \\
  \hline
  6 & Tentar adicionar 8\textordmasculine{} membro & ID do 8\textordmasculine{} usuário & ``Equipe já atingiu o limite de 7 membros.''; não adiciona. \\
  \hline
\end{longtable}

\subsection*{Caso de teste \#7 – Progresso de missão diária}

\textbf{Propósito:} Verificar se o progresso da missão diária é atualizado corretamente e se o XP é concedido ao concluir.

\textbf{Pré-condições:}
\begin{itemize}
  \item Usuário logado com missão diária ativa (meta: 1 reclamação).
\end{itemize}

\begin{longtable}{|c|p{4.5cm}|p{3cm}|p{4cm}|}
  \hline
  \textbf{\#} & \textbf{Ação} & \textbf{Estado anterior} & \textbf{Saída esperada} \\
  \hline
  \endfirsthead
  \hline
  \textbf{\#} & \textbf{Ação} & \textbf{Estado anterior} & \textbf{Saída esperada} \\
  \hline
  \endhead
  1 & Usuário cria uma reclamação dentro do período da missão & Progresso: 0/1 & Progresso atualizado para 1/1; missão concluída; XP adicionado; notificação \texttt{mission\_completed}. \\
  \hline
  2 & Usuário cria reclamação após a expiração da missão & Missão expirada & Progresso não é contabilizado; missão permanece não concluída. \\
  \hline
  3 & Usuário abre a aba de missões sem missão atribuída & Sem missão no período & Backend sorteia e atribui nova missão automaticamente; tela exibe a missão. \\
  \hline
\end{longtable}

\subsection*{Caso de teste \#9 – Alternância de tema (modo claro/escuro)}

\textbf{Propósito:} Verificar se o usuário consegue alternar entre modo claro
e modo escuro pelas configurações e se o tema é aplicado em toda a interface.

\textbf{Pré-condições:}
\begin{itemize}
  \item Usuário logado no aplicativo.
\end{itemize}

\begin{longtable}{|c|p{4cm}|p{4cm}|p{4cm}|}
  \hline
  \textbf{\#} & \textbf{Ação} & \textbf{Estado inicial} & \textbf{Saída esperada} \\
  \hline
  \endfirsthead
  \hline
  \textbf{\#} & \textbf{Ação} & \textbf{Estado inicial} & \textbf{Saída esperada} \\
  \hline
  \endhead
  1 & Ativar modo escuro nas configurações & Modo claro ativo & Interface muda para modo escuro em todas as telas. \\
  \hline
  2 & Ativar modo claro nas configurações & Modo escuro ativo & Interface muda para modo claro em todas as telas. \\
  \hline
  3 & Fechar e reabrir o aplicativo & Modo escuro salvo & Aplicativo reabre mantendo o modo escuro. \\
  \hline
\end{longtable}
